\documentclass[10pt,a4paper]{amsart}
\usepackage[margin=19mm]{geometry}
\usepackage{amsmath,amssymb,mathtools}
\usepackage[scr=rsfs,cal=euler]{mathalfa}
\usepackage{xfrac}
\usepackage[unicode,hidelinks,bookmarks=false]{hyperref}
\newcommand{\ie}{i.e.\ }
\newcommand{\cf}{cf.\ }
\newcommand{\resp}{respectively\ }
\newcommand{\Subsection}{\subsection}
\providecommand{\nobreakdash}{\nobreak\hbox{-}}
\setlength{\emergencystretch}{2em}
\multlinegap0pt
\usepackage{amssymb}
\newtheorem{theorem}{Theorem}
\newcommand{\R}{{\mathbb{R}}}
\newcommand{\mh}{\textsc{m}}
\title{Homogeneity\\}
\author{Martin Himmel}
\date{Source: arXiv:2509.25227v2 (CC BY 4.0)}
\begin{document}
\maketitle

\noindent\emph{Source and licence.} Homogeneity\\. Version arXiv:2509.25227v2 (CC BY 4.0), distributed by arXiv under the Creative Commons Attribution 4.0 International licence, http://creativecommons.org/licenses/by/4.0/, as recorded in the arXiv licence metadata for this exact version and shown on the abstract page https://arxiv.org/abs/2509.25227v2. Full licence notice: \texttt{licenses/arxiv-2509.25227-CC-BY-4.0.txt}.\par\smallskip

\noindent\textit{Editorial scope.} Counted unit: the article's theorem eq:multHomIndx (source0.tex lines 2458--2475). The article's complete original proof of it is delivered with it (source0.tex lines 2477--2530, one \texttt{proof} environment of 54 lines). No mathematics is written, repaired or re-derived: the statement and the proof are the article's own lines, in the article's own notation, and no retained line is substituted or re-flowed.  No further source-local context is needed: every cross-reference of every retained line resolves inside the two delivered spans.  Nothing was omitted from any retained span, so the package carries no omission record.  No retained line contains a citation, so the wrapper carries no bibliography and no bibliography entry.  No retained line contains a figure environment, an image-inclusion command or drawing code, so no image is delivered and the package declares no asset dependency.  Editorial formatting only, in addition: an AMS \texttt{amsart} preamble that restates the article's own declarations and macros used by the retained lines, namely the environments \texttt{theorem} on separate counters, together with the standard packages those lines need.  The retained environments print in this document as Theorem 1 in order of appearance, whereas the article prints the same items with the numbering of its own full document.  The complete native source of the article as distributed in the arXiv e-print for this exact version is bundled unchanged as \texttt{sources/186005/source0.tex}.
\begin{theorem}
Let $I \subset \R$ be an interval and $f:I \to (0, +\infty)$ be a differentiable function.
Then the multiplicative homogeneity function of $f$ is independent of the variable
\begin{enumerate}
	\item 
	$x$, if and only if
	\begin{equation*}
	\mh_f(x,t) = t \mh_{f^\prime} (x,t)
	\label{eq:multHomIndx}
	\end{equation*}
	for all $x \in I, t>0$,
    implying that $f$ is constant.
	
	\item 
	$t$, if, and only if, $f^\prime \equiv 0$ on $I$,
     whence $f$ is constant.
\end{enumerate}
\end{theorem}

\begin{proof} 
%(Sketch)
\begin{enumerate}
    \item 
    Since $f$ is differentiable and $\mh_f$ independent of $x$, we have
    \begin{equation*}
\frac{\partial}{\partial x} \mh_f(x,t) = 
\frac{t f^\prime (tx) f(x) - f(tx) f^\prime (x)}{(f(x))^2}
=0,
\end{equation*}
whence
\begin{equation*}
t f^\prime (tx) f(x) = f(tx) f^\prime (x),
\end{equation*}
Rearranging gives us
\begin{equation*}
t \frac{f^\prime (tx)}{f(tx)}  =
\frac{f^\prime (x)}{f(x)},
\end{equation*}
which is condition \eqref{eq:multHomIndx}.
Thus
\begin{equation*}
(\log(f(tx)))^\prime = \log(f (x))^\prime.
\end{equation*}
Integrating here both sides, we get
\begin{equation*}
\log(f(tx))) = \log(f(x)) + C(t),
\end{equation*}
where $C(t)$ is an arbitrary constant.
Thus,
\begin{equation*}
f(tx)) = f(x)  e^{C(t)},
\end{equation*}
For $t=1$ it follows $f(x)=\frac{f(1)}{e^{C(1)}}$, as claimed.
The conversion holds true,
since the only functions obeying \eqref{eq:multHomIndx} are constant, which are clearly independent of any variable $x$.

  \item 
Assume $\mh_f$ independent of $t$.
Since $f$ is differentiable we have
    \begin{equation*}
\frac{\partial}{\partial t} \mh_f(x,t) = 
\frac{x f^\prime (tx) - f(tx) \cdot 0}{(f(x))^2}
=0.
	%\label{eq:}
\end{equation*}
Thus, $x f^\prime (tx) =0$.
For $x=0$ this clearly holds;
assuming $x \neq 0$ 
we get $f^\prime (tx) =0$,
which means that $f$ is constant on an interval.
%The rest follows from the definition of the homogeneity function.
\end{enumerate} 
\end{proof}

\end{document}
